\section{Stacking the system into one tile}
\label{sec:stacking}

We now replace the finite system of \cref{model:system} by one
translational tiling equation.  The extra finite coordinate must preserve
cyclicity, so we use a fresh prime cyclic group.  The argument is a
version of the stacking construction in
\citet[Theorem~3.1 and Lemma~3.2]{GTPeriodic}, with earlier
consolidation results in \citet[Theorem~1.15 and Section~3]{GTTwoTiles}.
We give the partition and uniqueness arguments in full, choosing a
fresh prime to retain the cyclic finite factor.

\begin{lemma}[A partition with full difference sets]
\label{stack:partition}
For every positive integer $s$ and every finite set of primes, there is
a prime $q$ outside that set and a partition
\[
  \Zmod{q}=E_1\sqcup\cdots\sqcup E_s
\]
into nonempty sets satisfying $E_\nu-E_\nu=\Zmod{q}$ for every $\nu$.
\end{lemma}

\begin{proof}
Choose a sufficiently large odd prime $q$, and color the elements of
$\Zmod{q}$ independently and uniformly with $s$ colors.  Fix a nonzero
$d\in\Zmod{q}$ and a color $\nu$.  Since $q$ is prime, the elements
$0,d,2d,\ldots,(q-1)d$ form one cycle.  The pairs
\[
 (0,d),\ (2d,3d),\ldots,
 ((q-3)d,(q-2)d)
\]
are disjoint, and the second element minus the first is $d$ in every
pair.  The probability that a given pair has color $\nu$ at both ends
is $s^{-2}$; these events are independent across the pairs.  Therefore
the probability that $d\notin E_\nu-E_\nu$ is at most
\[
 (1-s^{-2})^{\lfloor q/2\rfloor}.
\]
By the union bound, the probability that some color fails to realize
some nonzero difference is at most
\[
 s(q-1)(1-s^{-2})^{\lfloor q/2\rfloor}.
\]
For $s=1$ this is zero, and for fixed $s>1$ it tends to zero as $q$
tends to infinity.  Choose $q$ large enough that the bound is less than
one, also avoiding the specified finite set of primes.  A successful
coloring gives every nonzero difference in every color.  Each color
class is consequently nonempty, which also gives the zero difference.
\end{proof}

\begin{proposition}[Stacking and periodicity]
\label{stack:equivalence}
Let $H$ be an abelian group, let $F_1,\ldots,F_s\subset H$ be finite
nonempty sets, and let $E_1,\ldots,E_s$ be a partition as in
\cref{stack:partition}.  Define
\[
 F^{\mathrm{st}}=\bigcup_{\nu=1}^s(F_\nu\times E_\nu)
 \subset H\times\Zmod{q}.
\]
There is a common solution to $A\oplus F_\nu=H$ for all $\nu$ if and
only if $F^{\mathrm{st}}$ tiles $H\times\Zmod{q}$.  Moreover, a
fully periodic tiling by $F^{\mathrm{st}}$ gives a fully periodic
common solution, and a fully periodic common solution gives a fully
periodic tiling by $F^{\mathrm{st}}$.
\end{proposition}

\begin{proof}
If $A$ is a common solution, then $A\times\{0\}$ is a tiling set for
$F^{\mathrm{st}}$: for a target $(g,t)$, the partition of $\Zmod{q}$
selects a unique $\nu$ with $t\in E_\nu$, and the equation
$A\oplus F_\nu=H$ selects a unique $g=a+f$.  If $A$ has a finite-index
period subgroup $P\leq H$, then $P\times\{0\}$ has finite index in
$H\times\Zmod{q}$ and preserves this tiling set.

Conversely, suppose that $B\oplus F^{\mathrm{st}}=H\times\Zmod{q}$,
and fix $g\in H$.  A contribution to the fibre $\{g\}\times\Zmod{q}$
of color $\nu$ consists of a pair $((a,t),f)$ with $(a,t)\in B$,
$f\in F_\nu$, and $a+f=g$.  This contribution covers the vertical
set $t+E_\nu$.  There are only finitely many contributions: $a$ is
determined by $f$, and there are at most $q$ possible values of $t$.
Two distinct contributions of the same color cannot occur.  Indeed,
the identity $E_\nu-E_\nu=\Zmod{q}$ makes any two translates of
$E_\nu$ intersect, contradicting uniqueness in the stacked tiling.

Let $n_\nu(g)\in\{0,1\}$ be the number of contributions of color
$\nu$.  Exact coverage of the fibre gives
\begin{equation}
 \sum_{\nu=1}^s n_\nu(g)|E_\nu|
   =q=\sum_{\nu=1}^s|E_\nu|.
 \label{stack:fibre-count}
\end{equation}
Every $|E_\nu|$ is positive, so \cref{stack:fibre-count} forces
$n_\nu(g)=1$ for all $\nu$.

Projection $\operatorname{pr}:H\times\Zmod{q}\to H$ is injective on
$B$.  Otherwise, two distinct points $(a,t),(a,t')\in B$, together
with any $f\in F_1$, would give two color-$1$ contributions at $g=a+f$.
Consequently $A=\operatorname{pr}(B)$ is a set of projected
translations without multiplicity.  The equality $n_\nu(g)=1$ now
says exactly that $g$ has one representation in $A+F_\nu$.  Thus
$A\oplus F_\nu=H$ for every $\nu$.

Finally, if $B+P=B$ for a finite-index subgroup
$P\leq H\times\Zmod{q}$, then
$A+\operatorname{pr}(P)=A$.  The induced map
\[
 (H\times\Zmod{q})/P\longrightarrow H/\operatorname{pr}(P)
\]
is surjective, so $\operatorname{pr}(P)$ has finite index.  This proves
the periodicity assertion.
\end{proof}

\begin{corollary}[One tile with a cyclic finite coordinate]
\label{stack:cyclic}
For some positive integer $Q$, there is a finite nonempty set
$F\subset\Gamma=\Z^2\times\Zmod{Q}$ which tiles $\Gamma$ but admits
no fully periodic tiling.
\end{corollary}

\begin{proof}
Apply \cref{stack:equivalence} to the finite system from
\cref{model:system} in $G=\Z^2\times V$.  Choose the prime $q$ in
\cref{stack:partition} not to divide $|V|$.  Since $V$ is cyclic,
$V\times\Zmod{q}$ is cyclic of order $Q=q|V|$: a pair of generators
has order $Q$ by coprimality.  The common solution gives a tiling by
the stacked tile, and a fully periodic stacked tiling would give a
fully periodic common solution, contrary to \cref{model:system}.
Identify the finite product with $\Zmod{Q}$ and call the stacked tile
$F$. Since every system tile has cardinality $M=|H_0|$, the
disjoint stacking gives $|F|=qM$.
\end{proof}
